\chapter*{ABSTRACT}
\addcontentsline{toc}{chapter}{ABSTRACT}
\bigskip

El Ni\~{n}o events can substantially influence the South Asian monsoon, creating important implications for agriculture, water resources, and disaster preparedness. This project develops a deep learning framework for forecasting the El Ni\~{n}o--Southern Oscillation (ENSO) state 3 to 6 months in advance and assessing its potential impacts on South Asian precipitation and temperature. The framework uses climate data from 1980 to 2025, incorporating publicly available reanalysis datasets, including ERA5 and ORAS5, together with established ENSO indices. A hybrid CNN-TCN architecture is employed, where the CNN component extracts spatial patterns from monthly sea surface temperature (SST) fields, while the TCN component captures their temporal evolution over the 12-month input sequence. The primary prediction target is the Ni\~{n}o 3.4 index, which is used as an indicator of the ENSO state. Evaluated using a 5-seed ensemble, the model achieves a mean 3-month lead Ni\~{n}o 3.4 RMSE of 0.51\degree C and a Pearson correlation of 0.92. The predicted ENSO state is subsequently used to derive an ENSO-based assessment of South Asian precipitation and temperature anomalies. The resulting pipeline provides a deterministic ENSO-based outlook and corresponding impact assessment for the 2026 June--September monsoon season, demonstrating the potential of deep learning-based ENSO forecasting to support climate-risk assessment, disaster preparedness, and resource management in South Asia.

\vspace{1em}
\noindent\textit{\textbf{Keywords}}: \textit{CNN-TCN, Deep Learning, El Ni\~{n}o, ENSO, Ni\~{n}o 3.4, Sea Surface Temperature}.

\clearpage